\Nguide{} steht weiterhin im \emph{Leitfaden für Folgen}.
Hier stehen seine separaten Beweistabellen; die
\href{_B21-newton-reading.pdf\#newton.reading}{Lesefassung}
erklärt die Ideen und Beispiele. Die Hilfssätze erhalten ausschließlich
hier die Kennung $21N$. Kurze Buchstabenverweise führen jeweils zur
vollständigen Aussage und Tabelle.

\section*{Notation und verwendete Grundlagen}
Alle Folgen sind reell und beginnen bei $0$.
Soweit im jeweiligen Satz nicht ausdrücklich anders typisiert, gelten die Typen $a,b,u,v\in\Nseq$, $n,m,k,r,s,d,j\in\mathbb N$
und $\alpha,\beta\in\mathbb R$ als ständiger Kontext.
Die linke Spalte nennt die zusätzlich offenen Annahmen.
Alle Aussagen über Folgen sind gliedweise zu lesen; gebundene
Allverallgemeinerung entlässt die jeweilige bloße Typannahme.
Wir verwenden die Differenzenrekursion aus
\hyperref[b21guide:def:differences]{Definition~\ref*{b21guide:def:differences}}:
\[
 (\Delta^0a)_n=a_n,\qquad
 (\Delta^{r+1}a)_n=(\Delta^ra)_{n+1}-(\Delta^ra)_n.
\]
Die Summen- und Potenznotation wird rekursiv festgelegt durch
\[
 \sigma_0(u)=0,\quad \sigma_{m+1}(u)=\sigma_m(u)+u_m,
 \qquad x^0=1,\quad x^{j+1}=x^jx.
\]
Dabei ist $\sigma_m(u)=\sum_{i<m}u_i$ und insbesondere $0^0=1$.
Leere Summen sind $0$, leere Produkte $1$; Produkte werden entsprechend
rekursiv gebildet. Wir setzen
\[
 \begin{gathered}
 X_n=n,\quad u_j(n)=n^j,\quad D_k(a)=(\Delta^ka)_0,\\
 P_c=\sum_{i=0}^d c_i u_i\quad(c\in\mathbb R^{d+1}),\qquad
 \mathcal P_d=\{P_c:c\in\mathbb R^{d+1}\},\quad\mathcal P_{-1}=\{0\},\\
 B_0=1,\qquad B_k(n)=\frac{\prod_{i=0}^{k-1}(n-i)}{k!}\quad(k\geq1).
 \end{gathered}
\]
$0$ und $1$ bezeichnen bei Folgen die konstanten Folgen.
$\mathcal P_d$ ist zunächst nur eine Menge dargestellter Folgen;
Eindeutigkeit der Koeffizienten wird nicht vorausgesetzt.
Außerhalb eines angegebenen Koeffiziententupels werden dessen
Koeffizienten mit $0$ fortgesetzt. Für natürliche $n$ gilt die
Binomialkonvention des Leitfadens; ergänzend ist $\binom n{-1}=0$.

Die kurzen Begründungen verwenden folgende bereits vorhandene Grundlagen:
\begin{itemize}
\item $\NR$: Körperrechnung (Assoziativität, Kommutativität,
Distributivität, neutrale Elemente, additive Inverse und Kürzen mit
nachweislich von $0$ verschiedenen Zahlen) aus
\FormulaRefAuto{RealCutsCompleteOrderedField}[thm] und
\FormulaRefAuto{CompleteOrderedFieldAxioms}[def]. Die Einbettung der
natürlichen Zahlen wird wie im Leitfaden stillschweigend verwendet.
\item $\NF$: \FormulaRefAuto{FactorialRecursion}[thm] und
\FormulaRefAuto{FactorialNonzero}[thm].
\item $\NI{\cdot}$: \FormulaRefAuto{PeanoInductionRuleAddOne}[thm].
Die angegebenen Zeilen bezeichnen Anfang und Induktionsschritt.
\end{itemize}
Ein mit $\NR$ bezeichneter Schritt ist eine endliche algebraische
Umformung, kein zusätzlich angenommener Folgensatz. Die benötigten
Summen-, Binomial- und Differenzenregeln werden im Folgenden bewiesen.
Die Grundlagen werden einmal hier genannt, damit sie nicht nach jeder
Tabelle erneut ausgeschrieben werden müssen.

% Lokale Grundlagen für die Newton-Beweistabellen; keine Leitfadenbeweise.
\section*{Endliche Rechnung und Differenzen}

\Nlemma{S}{Reelle endliche Summen}{%
 \begin{gathered}
 \sigma_m(\alpha u+\beta v)=\alpha\sigma_m(u)+\beta\sigma_m(v),
 \qquad \sigma_m(0)=0,\\
 \sigma_{m+1}(u)=u_0+\sigma_m(u^+),\qquad u^+_i=u_{i+1},\\
 \bigl(\forall j<r\;u_{m+j}=0\bigr)
 \ \Longrightarrow\ \sigma_{m+r}(u)=\sigma_m(u).
 \end{gathered}
}{NewtonFiniteSums}
Für die folgende Tabelle seien $L_m$ die erste Gleichung, $T_m$ die
Verschiebungsgleichung und
$Z_r=(\forall j<r\;u_{m+j}=0)\to\sigma_{m+r}(u)=\sigma_m(u)$.
Die Parameter der jeweiligen Induktion sind beliebig festgehalten.
\begin{tabproofwide}
 \proofstepwidestar[]{L_0}{\Nwhy{Summendefinition; \(0=\alpha0+\beta0\)}}
 \proofstepwidestar[2]{L_m}{\rA}
 \proofstepwidestar[2]{\begin{aligned}[t]
 \sigma_{m+1}(\alpha u+\beta v)
 &=\sigma_m(\alpha u+\beta v)+\alpha u_m+\beta v_m\\
 &=\alpha(\sigma_m(u)+u_m)+\beta(\sigma_m(v)+v_m)
 \end{aligned}}{\Nwhy{Summendefinition; 2; Körperrechnung}}
 \proofstepwidestar[2]{L_{m+1}}{\Nwhy{3; Summendefinition}}
 \proofstepwidestar[]{\forall m\in\mathbb N\;L_m}{\NI{1;2--4}}
 \proofstepwidestar[]{\sigma_m(0)=0}{\Nwhy{5 mit \(\alpha=\beta=0\); Körperrechnung}}
 \proofstepwidestar[]{T_0}{\Nwhy{\(\sigma_1(u)=u_0=u_0+\sigma_0(u^+)\)}}
 \proofstepwidestar[8]{T_m}{\rA}
 \proofstepwidestar[8]{\begin{aligned}[t]
 \sigma_{m+2}(u)&=\sigma_{m+1}(u)+u_{m+1}\\
 &=u_0+\sigma_m(u^+)+u^+_m
 =u_0+\sigma_{m+1}(u^+)
 \end{aligned}}{\Nwhy{Summendefinition; 8; Def. \(u^+\)}}
 \proofstepwidestar[]{\forall m\in\mathbb N\;T_m}{\NI{7;8--9}}
 \proofstepwidestar[]{Z_0}{\Nwhy{\(\sigma_{m+0}(u)=\sigma_m(u)\); \(\to I\)}}
 \proofstepwidestar[12]{Z_r}{\rA}
 \proofstepwidestar[13]{\forall j<r+1\;u_{m+j}=0}{\rA}
 \proofstepwidestar[13]{(\forall j<r\;u_{m+j}=0)\land u_{m+r}=0}{\Nwhy{13; \(j<r\Rightarrow j<r+1\), \(r<r+1\)}}
 \proofstepwidestar[12,13]{\sigma_{m+r}(u)=\sigma_m(u)}{\rRE{12,\rAEa{14}}}
 \proofstepwidestar[12,13]{\sigma_{m+r+1}(u)=\sigma_{m+r}(u)+u_{m+r}=\sigma_m(u)}{\Nwhy{Summendefinition; 14,15}}
 \proofstepwidestar[12]{Z_{r+1}}{\rRI{13,16}}
 \proofstepwidestar[]{\forall r\in\mathbb N\;Z_r}{\NI{11;12--17}}
\end{tabproofwide}
Die Nullverlängerung gilt insbesondere für die Nullfolge. Mit $m=0$
besagt sie auch: Sind alle tatsächlich summierten Glieder null, so ist
die Summe null. Die Verschiebung ist die benötigte Indexumbenennung;
allgemeine Umordnungen unendlicher Summen werden nicht verwendet.

\Needspace{14\baselineskip}
\Nlemma{C}{Koeffizientenrechnung}{%
 \begin{gathered}
 \alpha P_c+\beta P_e=P_{\alpha c+\beta e},\qquad
 (\alpha X+\beta)P_c=P_b,\quad b_i=\alpha c_{i-1}+\beta c_i,\\
 d\le D\ \Longrightarrow\ \mathcal P_d\subseteq\mathcal P_D,\\
 p,q\in\mathcal P_d\ \Longrightarrow\
 \alpha p+\beta q\in\mathcal P_d\ \land\
 (\alpha X+\beta)p\in\mathcal P_{d+1},\\
 \forall i<m\;(p_i\in\mathcal P_d)
 \ \Longrightarrow\ \sum_{i<m}p_i\in\mathcal P_d.
 \end{gathered}
}{NewtonCoefficientCalculus}
In der ersten Zeile sind $c,e\in\mathbb R^{d+1}$,
$b\in\mathbb R^{d+2}$ und $c_i=0$ außerhalb von $0\le i\le d$.
Die Tupelargumente der Tabelle gelten für $d\ge0$; die eigenen Nullfolgenzeilen behandeln $d=-1$.
Bei $D\ge d$ bezeichnet $\widetilde c$ die bis $D$ mit Nullen
verlängerte Koeffizientenliste. Die letzte Aussage wird mit
$F_m=(\forall i<m\;p_i\in\mathcal P_d)\to\sum_{i<m}p_i\in\mathcal P_d$
abgekürzt.
\begin{tabproofwide}
 \proofstepwidestar[]{\alpha P_c+\beta P_e
 =\sum_{i=0}^d(\alpha c_i+\beta e_i)u_i=P_{\alpha c+\beta e}}
 {\Nwhy{Def. \(P_c\); \Nref{S}; punktweise Auswertung}}
 \proofstepwidestar[]{\begin{aligned}[t]
 (\alpha X+\beta)P_c
 &=\alpha\sum_{i=0}^d c_i u_{i+1}
   +\beta\sum_{i=0}^d c_i u_i\\
 &=\sum_{i=0}^{d+1}(\alpha c_{i-1}+\beta c_i)u_i=P_b
 \end{aligned}}{\Nwhy{\Nref{S}; \(Xu_i=u_{i+1}\); Nullverlängerung}}
 \proofstepwidestar[]{0\in\mathcal P_d}{\Nwhy{Nullkoeffizienten; \Nref{S}; für \(d=-1\) Definition}}
 \proofstepwidestar[4]{p,q\in\mathcal P_d}{\rA}
 \proofstepwidestar[4]{\exists c,e\in\mathbb R^{d+1}\;(p=P_c\land q=P_e)}{\Nwhy{Def. \(\mathcal P_d\), zunächst \(d\ge0\)}}
 \proofstepwidestar[6]{c,e\in\mathbb R^{d+1}\land p=P_c\land q=P_e}{\rA}
 \proofstepwidestar[6]{\alpha p+\beta q\in\mathcal P_d}{\Nwhy{1; 6; Def. \(\mathcal P_d\)}}
 \proofstepwidestar[6]{(\alpha X+\beta)p\in\mathcal P_{d+1}}{\Nwhy{2; 6; Def. \(\mathcal P_{d+1}\)}}
 \proofstepwidestar[6]{\alpha p+\beta q\in\mathcal P_d\land
 (\alpha X+\beta)p\in\mathcal P_{d+1}}{\rAI{7,8}}
 \proofstepwidestar[4]{\alpha p+\beta q\in\mathcal P_d\land
 (\alpha X+\beta)p\in\mathcal P_{d+1}}{\rEE{5,6,9}}
 \proofstepwidestar[]{p,q\in\mathcal P_d\to
 \alpha p+\beta q\in\mathcal P_d\land
 (\alpha X+\beta)p\in\mathcal P_{d+1}}{\rRI{4,10}}
 \proofstepwidestar[]{\begin{aligned}[t]
 p,q\in\mathcal P_{-1}\to{}&\alpha p+\beta q\in\mathcal P_{-1}\\[-2pt]
 &\land\ (\alpha X+\beta)p\in\mathcal P_0
 \end{aligned}}{\Nwhy{\(p=q=0\); 3; Körperrechnung}}
 \proofstepwidestar[13]{d\le D\land p\in\mathcal P_d}{\rA}
 \proofstepwidestar[13]{\exists c\in\mathbb R^{d+1}\;(p=P_c)}{\Nwhy{13; Def. \(\mathcal P_d\), zunächst \(d\ge0\)}}
 \proofstepwidestar[15]{c\in\mathbb R^{d+1}\land p=P_c}{\rA}
 \proofstepwidestar[13,15]{p\in\mathcal P_D}
 {\Nwhy{\Nref{S}, Nullverlängerung; 13,15; Def. \(\mathcal P_D\)}}
 \proofstepwidestar[13]{p\in\mathcal P_D}{\rEE{14,15,16}}
 \proofstepwidestar[]{d\le D\land p\in\mathcal P_d\to p\in\mathcal P_D}{\rRI{13,17}}
 \proofstepwidestar[]{\mathcal P_{-1}\subseteq\mathcal P_D}{\Nwhy{\(\mathcal P_{-1}=\{0\}\); 3, auch für \(D=-1\)}}
 \proofstepwidestar[]{F_0}{\Nwhy{Leere Summe; 3; \(\to I\)}}
 \proofstepwidestar[21]{F_m}{\rA}
 \proofstepwidestar[22]{\forall i<m+1\;(p_i\in\mathcal P_d)}{\rA}
 \proofstepwidestar[22]{(\forall i<m\;p_i\in\mathcal P_d)\land p_m\in\mathcal P_d}{\Nwhy{22; Indexeinschränkung}}
 \proofstepwidestar[21,22]{\sum_{i<m}p_i\in\mathcal P_d}{\rRE{21,\rAEa{23}}}
 \proofstepwidestar[21,22]{\sum_{i<m+1}p_i\in\mathcal P_d}{\Nwhy{Summenrekursion; 11--12 mit \(\alpha=\beta=1\); 23,24}}
 \proofstepwidestar[21]{F_{m+1}}{\rRI{22,25}}
 \proofstepwidestar[]{\forall m\in\mathbb N\;F_m}{\NI{20;21--26}}
\end{tabproofwide}
Die Koeffizientenlisten werden ausschließlich als Darstellungen benutzt;
ihre Eindeutigkeit ist hier keine Voraussetzung.

\Nlemma{P}{Pascal-Identität mit Randfällen}{%
 \binom{n+1}{k}=\binom nk+\binom n{k-1}
 \qquad(n,k\in\mathbb N),\qquad \binom n{-1}=0.
}{NewtonPascal}
\begin{tabproofwide}
 \proofstepwidestar[1]{1\le k\le n}{\rA}
 \proofstepwidestar[1]{\begin{aligned}[t]
 \binom nk+\binom n{k-1}
 &=\frac{n!(n-k+1)}{k!(n-k+1)!}
   +\frac{n!k}{k!(n-k+1)!}\\
 &=\frac{n!(n+1)}{k!(n+1-k)!}
 =\binom{n+1}k
 \end{aligned}}{\Nwhy{Binomialdefinition; \NF; Körperrechnung}}
 \proofstepwidestar[]{1\le k\le n\to
 \binom{n+1}k=\binom nk+\binom n{k-1}}
 {\Nwhy{\(\to I\), 1,2; Symmetrie}}
 \proofstepwidestar[]{\binom{n+1}0=1=\binom n0+\binom n{-1}}
 {\Nwhy{Definition; \NF; \(0!=1\)}}
 \proofstepwidestar[]{\binom{n+1}{n+1}=1=\binom n{n+1}+\binom nn}
 {\Nwhy{Definition; \NF}}
 \proofstepwidestar[6]{k>n+1}{\rA}
 \proofstepwidestar[6]{\binom{n+1}k=0=\binom nk+\binom n{k-1}}
 {\Nwhy{Definition; \(k>n\), \(k-1>n\)}}
 \proofstepwidestar[]{k>n+1\to
 \binom{n+1}k=\binom nk+\binom n{k-1}}{\rRI{6,7}}
 \proofstepwidestar[]{k=0\ \lor\ 1\le k\le n\ \lor\ k=n+1\ \lor\ k>n+1}
 {\Nwhy{Natürliche Ordnung; vollständige Fälle}}
 \proofstepwidestar[]{\binom{n+1}k=\binom nk+\binom n{k-1}}
 {\Nwhy{3--5,8,9; Fallunterscheidung}}
\end{tabproofwide}
Für $n=0$ ist der innere Fall leer; die Zeilen 4, 5 und 8 liefern
unverändert sämtliche Werte von $k$.

\Nlemma{D}{Lineare Differenzen und Iteration}{%
 \begin{gathered}
 \Delta^r(\alpha a+\beta b)=\alpha\Delta^ra+\beta\Delta^rb,
 \qquad \Delta^r0=0,\\
 \Delta^{r+s}a=\Delta^r(\Delta^sa),\\
 \Delta^r\left(\sum_{i<m}\alpha_i a^{(i)}\right)
 =\sum_{i<m}\alpha_i\Delta^ra^{(i)}.
 \end{gathered}
}{NewtonDifferenceCalculus}
Die endliche Familie erfüllt $\alpha_i\in\mathbb R$ und $a^{(i)}\in\mathbb R^{\mathbb N}$ für $i<m$. Seien $L_r$ die Linearitätsgleichung, $K_r$ die Kompositionsgleichung
bei festem $s$ und $M_m$ die Gleichung für die endliche Summe.
\begin{tabproofwide}
 \proofstepwidestar[]{\begin{aligned}[t]
 (\Delta(\alpha a+\beta b))_n
 &=(\alpha a_{n+1}+\beta b_{n+1})-(\alpha a_n+\beta b_n)\\
 &=\alpha(\Delta a)_n+\beta(\Delta b)_n
 \end{aligned}}{\Nwhy{Differenzendefinition; Körperrechnung}}
 \proofstepwidestar[]{\Delta(\alpha a+\beta b)=\alpha\Delta a+\beta\Delta b}
 {\Nwhy{1; Funktionenextensionalität}}
 \proofstepwidestar[]{L_0}{\Nwhy{\(\Delta^0a=a\)}}
 \proofstepwidestar[4]{L_r}{\rA}
 \proofstepwidestar[4]{\begin{aligned}[t]
 \Delta^{r+1}(\alpha a+\beta b)
 &=\Delta(\alpha\Delta^ra+\beta\Delta^rb)\\
 &=\alpha\Delta^{r+1}a+\beta\Delta^{r+1}b
 \end{aligned}}{\Nwhy{Differenzenrekursion; 2,4}}
 \proofstepwidestar[]{\forall r\in\mathbb N\;L_r}{\NI{3;4--5}}
 \proofstepwidestar[]{\Delta^r0=0}{\Nwhy{6 mit \(\alpha=\beta=0\); Körperrechnung}}
 \proofstepwidestar[]{K_0}{\Nwhy{\(\Delta^{0+s}a=\Delta^sa=\Delta^0(\Delta^sa)\)}}
 \proofstepwidestar[9]{K_r}{\rA}
 \proofstepwidestar[9]{\Delta^{r+1+s}a
 =\Delta(\Delta^r(\Delta^sa))=\Delta^{r+1}(\Delta^sa)}
 {\Nwhy{Differenzenrekursion; 9; \((r+1)+s=(r+s)+1\)}}
 \proofstepwidestar[]{\forall r\in\mathbb N\;K_r}{\NI{8;9--10}}
 \proofstepwidestar[]{M_0}{\Nwhy{Leere Summe; 7}}
 \proofstepwidestar[13]{M_m}{\rA}
 \proofstepwidestar[13]{\begin{aligned}[t]
 \Delta^r\left(\sum_{i<m+1}\alpha_i a^{(i)}\right)
 &=\Delta^r\left(\sum_{i<m}\alpha_i a^{(i)}\right)
   +\alpha_m\Delta^ra^{(m)}\\
 &=\sum_{i<m+1}\alpha_i\Delta^ra^{(i)}
 \end{aligned}}{\Nwhy{Summenrekursion; 6,13}}
 \proofstepwidestar[]{\forall m\in\mathbb N\;M_m}{\NI{12;13--14}}
\end{tabproofwide}

\Nlemma{B}{Binomialfolgen sind polynomiale Folgen}{%
 B_k(n)=\binom nk,\qquad B_k\in\mathcal P_k,
 \qquad B_{k+1}=\frac{X-k}{k+1}B_k
 \quad(n,k\in\mathbb N).
}{NewtonBinomialPolynomials}
Setze $J_k=(\forall n\in\mathbb N\;B_k(n)=\binom nk)
\land B_k\in\mathcal P_k$.
\begin{tabproofwide}
 \proofstepwidestar[]{B_{k+1}=\frac{X-k}{k+1}B_k}
 {\Nwhy{Produktrekursion; \NF; Körperrechnung; Extensionalität}}
 \proofstepwidestar[2]{k<n}{\rA}
 \proofstepwidestar[2]{\begin{aligned}[t]
 \binom n{k+1}
 &=\frac{n!}{(k+1)!(n-k-1)!}\\
 &=\frac{n-k}{k+1}\frac{n!}{k!(n-k)!}
 =\frac{n-k}{k+1}\binom nk
 \end{aligned}}{\Nwhy{Binomialdefinition; \NF; Körperrechnung}}
 \proofstepwidestar[]{k<n\to\binom n{k+1}=\frac{n-k}{k+1}\binom nk}{\rRI{2,3}}
 \proofstepwidestar[]{k=n\to\binom n{k+1}=0=\frac{n-k}{k+1}\binom nk}
 {\Nwhy{Binomialdefinition; \(n-k=0\)}}
 \proofstepwidestar[]{k>n\to\binom n{k+1}=0=\frac{n-k}{k+1}\binom nk}
 {\Nwhy{Binomialdefinition; \(\binom nk=0\)}}
 \proofstepwidestar[]{\binom n{k+1}=\frac{n-k}{k+1}\binom nk}
 {\Nwhy{4--6; \(k<n\lor k=n\lor k>n\)}}
 \proofstepwidestar[]{J_0}
 {\Nwhy{\(B_0=1\in\mathcal P_0\), \(\binom n0=1\); \NF}}
 \proofstepwidestar[9]{J_k}{\rA}
 \proofstepwidestar[9]{\forall n\in\mathbb N\;
 B_{k+1}(n)=\frac{n-k}{k+1}B_k(n)=\binom n{k+1}}
 {\Nwhy{1,7,9; \(\forall E,\forall I\)}}
 \proofstepwidestar[9]{B_{k+1}\in\mathcal P_{k+1}}
 {\Nwhy{1,9; \Nref{C} mit \(\alpha=1/(k+1)\), \(\beta=-k/(k+1)\)}}
 \proofstepwidestar[9]{J_{k+1}}{\Nwhy{10,11; Transitivität; \(\land I\)}}
 \proofstepwidestar[]{\forall k\in\mathbb N\;J_k}{\NI{8;9--12}}
\end{tabproofwide}


\section*{Die allgemeine Newtondarstellung}
\Nlemma{N}{Newtonidentität}{%
 \forall n\in\mathbb N\;\forall a\in\Nseq\quad
 a_n=\sum_{k=0}^n\binom nkD_k(a)
}{NewtonFiniteIdentity}
Setze $\mathsf N(n)\equiv
\forall a\in\Nseq\;(a_n=\sum_{k=0}^n\binom nkD_k(a))$.
Im Schritt ist $a$ beliebig, und $D_k$ steht für $D_k(a)$.
\begin{tabproofwide}
 \proofstepwidestar[]{\mathsf N(0)}{\Nwhy{Definitionen: $a_0=1\cdot(\Delta^0a)_0$.}}
 \proofstepwidestar[2]{\mathsf N(n)}{\rA}
 \proofstepwidestar[2]{a_n=\sum_{k=0}^n\binom nkD_k}{\rUE{2}}
 \proofstepwidestar[2]{(\Delta a)_n=\sum_{k=0}^n\binom nkD_{k+1}}{\Nwhy{$\forall E(2)$ für $\Delta a$; \Nref{D}.}}
 \proofstepwidestar[2]{\begin{aligned}[t]a_{n+1}
 &=a_n+(\Delta a)_n\\
 &=\sum_{k=0}^n\binom nkD_k
   +\sum_{k=1}^{n+1}\binom n{k-1}D_k
 \end{aligned}}{\Nwhy{Differenzenrekursion, (3,4), \Nref{S}.}}
 \proofstepwidestar[2]{a_{n+1}=\sum_{k=0}^{n+1}
   \left(\binom nk+\binom n{k-1}\right)D_k}{\Nwhy{\Nref{S}, Nullverlängerung, (5).}}
 \proofstepwidestar[2]{a_{n+1}=\sum_{k=0}^{n+1}\binom{n+1}kD_k}{\Nwhy{\Nref{P}, (6).}}
 \proofstepwidestar[2]{\mathsf N(n+1)}{\rUI{7}}
 \proofstepwidestar[]{\forall n\in\mathbb N\;\mathsf N(n)}{\NI{1;2\text{--}8}}
\end{tabproofwide}
Die Induktionsannahme gilt ausdrücklich für \emph{alle} Folgen;
deshalb ist ihre Anwendung auf $\Delta a$ in Zeile~4 zulässig.

\section*{Gradsenkung ohne allgemeinen binomischen Lehrsatz}
\Nlemma{G}{Eine Differenz und ihr oberster Koeffizient}{%
 \begin{aligned}
 &c\in\mathbb R^{d+2}\ \Longrightarrow\
 \exists b\in\mathbb R^{d+1}\quad
   \Delta P_c=P_b\ \land\ b_d=(d+1)c_{d+1},\\
 &c\in\mathbb R\ \Longrightarrow\ \Delta(cu_0)=0.
 \end{aligned}
}{NewtonDegreeLowering}
Schreibe $q_j=\Delta u_j$ und
$\mathsf Q(j)\equiv\exists h\in\mathcal P_{j-1}\;
(q_{j+1}=(j+1)u_j+h)$.
Eine solche Gleichung gilt als Folgengleichheit an allen Indizes zugleich.
\begin{tabproofwide}
 \proofstepwidestar[]{q_0=0,\quad q_1=u_0,\quad\mathsf Q(0)}{\Nwhy{Potenz- und Differenzenrekursion; $h=0$.}}
 \proofstepwidestar[2]{\mathsf Q(j)}{\rA}
 \proofstepwidestar[3]{h\in\mathcal P_{j-1}\ \land\ q_{j+1}=(j+1)u_j+h}{\Nwhy{Annahme für $\exists E$ aus (2); $h$ frisch.}}
 \proofstepwidestar[]{q_{j+2}=(X+1)q_{j+1}+u_{j+1}}{\Nwhy{Potenzrekursion, $\NR$: $(n+1)[(n+1)^{j+1}-n^{j+1}]+n^{j+1}$.}}
 \proofstepwidestar[3]{q_{j+2}=(j+2)u_{j+1}+\bigl((j+1)u_j+(X+1)h\bigr)}{\Nwhy{(3,4), $\NR$.}}
 \proofstepwidestar[3]{(j+1)u_j+(X+1)h\in\mathcal P_j}{\Nwhy{\Nref{C}; auch $j=0$ mit $h=0$.}}
 \proofstepwidestar[3]{\mathsf Q(j+1)}{\Nwhy{$\exists I$ aus (5,6).}}
 \proofstepwidestar[2]{\mathsf Q(j+1)}{\rEE{2,3,7}}
 \proofstepwidestar[]{\forall j\in\mathbb N\;\mathsf Q(j)}{\NI{1;2\text{--}8}}
 \proofstepwidestar[]{\begin{aligned}[t]\Delta P_c
 &=\sum_{j=0}^{d+1}c_jq_j\\
 &=c_{d+1}q_{d+1}+\sum_{j=1}^{d}c_jq_j
 \end{aligned}}{\Nwhy{\Nref{D}, \Nref{S}; $q_0=0$ aus (1).}}
 \proofstepwidestar[]{\forall j\,(1\leq j\leq d\ \Rightarrow\ q_j\in\mathcal P_{d-1})}{\Nwhy{(9), \Nref{C}: $q_j\in\mathcal P_{j-1}\subseteq\mathcal P_{d-1}$.}}
 \proofstepwidestar[]{\sum_{j=1}^{d}c_jq_j\in\mathcal P_{d-1}}{\Nwhy{(11), endliche Abgeschlossenheit aus \Nref{C}; für $d=0$ die Nullsumme.}}
 \proofstepwidestar[]{\Delta P_c\in(d+1)c_{d+1}u_d+\mathcal P_{d-1}}{\Nwhy{(9) mit $j=d$, (10,12), \Nref{C}.}}
 \proofstepwidestar[]{\exists b\in\mathbb R^{d+1}\;
 (\Delta P_c=P_b\land b_d=(d+1)c_{d+1})}{\Nwhy{(13), Definition von $\mathcal P_{d-1}$; höchsten Koeffizienten anhängen.}}
 \proofstepwidestar[]{\Delta(cu_0)=0}{\Nwhy{Differenzenrekursion, $\NR$.}}
\end{tabproofwide}
In Zeile~13 ist $v+\mathcal P_{d-1}=\{v+h:h\in\mathcal P_{d-1}\}$.
Der Rest in (13) besitzt Koeffizienten, die für die ganze Folge gelten;
es werden keine vom Index $n$ abhängigen Koeffizienten gewählt.

\Nlemma{I}{Iterierte Differenz und Abbruch}{%
 \forall c\in\mathbb R^{d+1}\quad
 \Delta^dP_c=d!c_d\,u_0\quad\land\quad\Delta^{d+1}P_c=0
}{NewtonIteratedLeadingCoefficient}
Die angezeigte Aussage sei $\mathsf T(d)$; die Induktion quantifiziert
wieder gleichzeitig über sämtliche Koeffiziententupel.
\begin{tabproofwide}
 \proofstepwidestar[]{\mathsf T(0)}{\Nwhy{$P_{(c_0)}=c_0u_0$, \Nref{G}, $0!=1$.}}
 \proofstepwidestar[2]{\mathsf T(d)}{\rA}
 \proofstepwidestar[]{\exists b\in\mathbb R^{d+1}\;
 (\Delta P_c=P_b\land b_d=(d+1)c_{d+1})}{\Nwhy{\Nref{G}; $c\in\mathbb R^{d+2}$ beliebig.}}
 \proofstepwidestar[4]{b\in\mathbb R^{d+1}\land\Delta P_c=P_b\land b_d=(d+1)c_{d+1}}{\Nwhy{Annahme für $\exists E$ aus (3); $b$ frisch.}}
 \proofstepwidestar[2,4]{\Delta^dP_b=d!b_d\,u_0\quad\land\quad\Delta^{d+1}P_b=0}{\Nwhy{$\forall E(2)$ mit (4).}}
 \proofstepwidestar[2,4]{\begin{aligned}[t]
 \Delta^{d+1}P_c&=(d+1)!c_{d+1}\,u_0,\\
 \Delta^{d+2}P_c&=0
 \end{aligned}}{\Nwhy{(4,5), \Nref{D}, $\NF$, $\NR$.}}
 \proofstepwidestar[2]{\Delta^{d+1}P_c=(d+1)!c_{d+1}\,u_0\ \land\ \Delta^{d+2}P_c=0}{\rEE{3,4,6}}
 \proofstepwidestar[2]{\mathsf T(d+1)}{\rUI{7}}
 \proofstepwidestar[]{\forall d\in\mathbb N\;\mathsf T(d)}{\NI{1;2\text{--}8}}
\end{tabproofwide}
Eine Annahme $c_d\ne0$ wird nicht benötigt. Das schließt die Nullfolge
und Darstellungen mit führenden Nullkoeffizienten ein.

\section*{Die drei äquivalenten Aussagen des Theorems}
\hypertarget{newton.proofs.equivalence}{}
Für festes $d$ und $a$ stehen die Abkürzungen
\[
 \begin{aligned}
 A_d&\equiv(\Delta^{d+1}a=0),\\
 C_d&\equiv\left(a=\sum_{k=0}^{d}D_k(a)B_k\right),\\
 Z_d&\equiv\exists c\in\mathbb R^{d+1}\;(a=P_c).
 \end{aligned}
\]
Sie entsprechen den Aussagen (1), (3), (2) in \Nguide{}.
\begin{tabproofwide}
 \proofstepwidestar[1]{A_d}{\rA}
 \proofstepwidestar[1]{\forall r\in\mathbb N\;(\Delta^{d+1+r}a=0)}{\Nwhy{(1), \Nref{D}: $\Delta^r0=0$.}}
 \proofstepwidestar[1]{C_d}{\Nwhy{\Nref{N}, \Nref{B}, (2), \Nref{S}. Bei $n<d$ gilt $B_k(n)=0$ für $k>n$; bei $n\geq d$ verschwinden $D_k(a)$ für $k>d$.}}
 \proofstepwidestar[]{A_d\rightarrow C_d}{\rRI{1,3}}
 \proofstepwidestar[5]{C_d}{\rA}
 \proofstepwidestar[5]{a\in\mathcal P_d\quad\land\quad Z_d}{\Nwhy{(5), \Nref{B}, \Nref{C}; Definition von $\mathcal P_d$.}}
 \proofstepwidestar[]{C_d\rightarrow Z_d}{\rRI{5,\rAEb{6}}}
 \proofstepwidestar[8]{Z_d}{\rA}
 \proofstepwidestar[9]{z\in\mathbb R^{d+1}\land a=P_z}{\Nwhy{Annahme für $\exists E$ aus (8); $z$ frisch.}}
 \proofstepwidestar[9]{\Delta^{d+1}P_z=0}{\Nwhy{\Nref{I}, (9).}}
 \proofstepwidestar[9]{A_d}{\Nwhy{(9,10), $=E$.}}
 \proofstepwidestar[8]{A_d}{\rEE{8,9,11}}
 \proofstepwidestar[]{Z_d\rightarrow A_d}{\rRI{8,12}}
 \proofstepwidestar[]{(A_d\rightarrow Z_d)\land(Z_d\rightarrow C_d)}{\Nwhy{Implikationstransitivität aus (4,7) und (13,4), $\land I$: \FormulaRefAuto{P\rightarrow Q, Q\rightarrow R\vdash P \rightarrow R}[thm].}}
 \proofstepwidestar[]{(A_d\leftrightarrow Z_d)\land(Z_d\leftrightarrow C_d)}{\rAI{\rLRI{\rAEa{14},13},\rLRI{\rAEb{14},7}}}
\end{tabproofwide}
Der Hauptsatz wird hier bewiesen, aber nicht mit einer zweiten
Satznummer neu deklariert.

\section*{Eindeutigkeit und Bestimmung aus Anfangswerten}
\Nlemma{U}{Eindeutigkeit der Potenzkoeffizienten}{%
 c,e\in\mathbb R^{d+1},\ P_c=P_e
 \quad\Longrightarrow\quad\forall i\leq d\;(c_i=e_i)
}{NewtonPowerCoefficientUniqueness}
$\mathsf U(d)$ bezeichne diese Aussage für alle $c,e$;
$c^-$ und $e^-$ entstehen durch Weglassen des letzten Koeffizienten.
\begin{tabproofwide}
 \proofstepwidestar[]{\mathsf U(0)}{\Nwhy{$P_{(c_0)}(0)=c_0$, Auswertung der Gleichheit bei $0$.}}
 \proofstepwidestar[2]{\mathsf U(d)}{\rA}
 \proofstepwidestar[3]{P_c=P_e\qquad(c,e\in\mathbb R^{d+2})}{\rA}
 \proofstepwidestar[3]{(d+1)!c_{d+1}\,u_0=(d+1)!e_{d+1}\,u_0}{\Nwhy{(3), \Nref{I}.}}
 \proofstepwidestar[3]{c_{d+1}=e_{d+1}}{\Nwhy{(4) bei $0$; $(d+1)!\ne0$ aus $\NF$; $\NR$.}}
 \proofstepwidestar[3]{P_{c^-}=P_{e^-}}{\Nwhy{(3,5), Summenrekursion; gleichen letzten Term abziehen.}}
 \proofstepwidestar[2,3]{\forall i\leq d+1\;(c_i=e_i)}{\Nwhy{(2,6) für $i\leq d$; (5) für $i=d+1$.}}
 \proofstepwidestar[2]{\mathsf U(d+1)}{\Nwhy{$\rightarrow I(3,7)$, $\forall I$ über $c,e$.}}
 \proofstepwidestar[]{\forall d\in\mathbb N\;\mathsf U(d)}{\NI{1;2\text{--}8}}
\end{tabproofwide}

\Nlemma{L}{Eine Differenz braucht nur benachbarte Werte}{%
 \bigl[\forall i\leq m\;(a_i=b_i)\bigr]\ \Longrightarrow\
 \bigl[\forall k,n\in\mathbb N\;(n+k\leq m\ \Rightarrow\
 (\Delta^ka)_n=(\Delta^kb)_n)\bigr]
}{NewtonDifferenceLocality}
Für festes $m,a,b$ sei
$\mathsf L(k)\equiv\forall n\in\mathbb N\;
(n+k\leq m\Rightarrow(\Delta^ka)_n=(\Delta^kb)_n)$.
\begin{tabproofwide}
 \proofstepwidestar[1]{\forall i\leq m\;(a_i=b_i)}{\rA}
 \proofstepwidestar[1]{\mathsf L(0)}{\Nwhy{(1), Definition von $\Delta^0$.}}
 \proofstepwidestar[3]{\mathsf L(k)}{\rA}
 \proofstepwidestar[4]{n+k+1\leq m}{\rA}
 \proofstepwidestar[3,4]{\begin{aligned}[t]
 (\Delta^ka)_{n+1}&=(\Delta^kb)_{n+1},\\
 (\Delta^ka)_n&=(\Delta^kb)_n
 \end{aligned}}{\Nwhy{(3,4); $n+k\leq m$ und $(n+1)+k\leq m$.}}
 \proofstepwidestar[3,4]{(\Delta^{k+1}a)_n=(\Delta^{k+1}b)_n}{\Nwhy{(5), Differenzenrekursion.}}
 \proofstepwidestar[3]{\mathsf L(k+1)}{\Nwhy{$\rightarrow I(4,6)$, $\forall I$.}}
 \proofstepwidestar[1]{\forall k\in\mathbb N\;\mathsf L(k)}{\NI{2;3\text{--}7}}
 \proofstepwidestar[]{\bigl[\forall i\leq m\;(a_i=b_i)\bigr]\rightarrow\forall k\in\mathbb N\;\mathsf L(k)}{\rRI{1,8}}
\end{tabproofwide}

\Needspace{15\baselineskip}
\noindent\textbf{Anfangswerte und der Randfall $d=0$.}
Die folgende Tabelle beendet die Zusatzbehauptungen von \Nguide{}.
\begin{tabproofwide}
 \proofstepwidestar[1]{\Delta^{d+1}a=0\ \land\ \Delta^{d+1}b=0}{\rA}
 \proofstepwidestar[2]{\forall i\leq d\;(a_i=b_i)}{\rA}
 \proofstepwidestar[2]{\forall k\leq d\;(D_k(a)=D_k(b))}{\Nwhy{\Nref{L} mit $m=d$, $n=0$.}}
 \proofstepwidestar[1,2]{a=\sum_{k=0}^dD_k(a)B_k
   =\sum_{k=0}^dD_k(b)B_k=b}{\Nwhy{Bewiesene Äquivalenz $A_d\Rightarrow C_d$, (1,3).}}
 \proofstepwidestar[]{\begin{aligned}[t]
 &\bigl[\Delta^{d+1}a=0\land\Delta^{d+1}b=0\bigr]\\
 &\quad\rightarrow\bigl([\forall i\leq d\;(a_i=b_i)]\rightarrow a=b\bigr)
 \end{aligned}}{\rRI{1,\rRI{2,4}}}
 \proofstepwidestar[]{\Delta a=0\ \leftrightarrow\ a=a_0u_0}{\Nwhy{Bewiesene Äquivalenz mit $d=0$; $B_0=u_0$ und $D_0(a)=a_0$.}}
\end{tabproofwide}
